\documentclass[12pt]{article}
\usepackage[utf8]{inputenc}
\usepackage[T1]{fontenc}
\usepackage{hyperref}
\title{Turung Language Resources for K-12 Curriculum Development}
\author{Compiled by Claude Sonnet 5 (claude-sonnet-5) for Andr\'as Kornai}
\date{}
\begin{document}
\maketitle
\section{Introduction}
Turung (ISO 639-3 try) is spoken in a few villages in Assam, India. Sources disagree on its classification. Wikipedia calls it Sino-Tibetan and close to Singpho, while Ethnologue-derived sources list it as a Tai language. The Turung people's own accounts say they were originally Tai speakers who adopted Singpho in the late 18th or early 19th century.
\section{Primary Educational Resources}
No school materials were found.
\section{Linguistic and Reference Materials}
\begin{itemize}
\item Overview: \url{https://en.wikipedia.org/wiki/Turung_language}
\item Documentation project of Turung and Singpho in Assam (Stephen Morey and colleagues): \url{https://cultureincrisis.org/projects/a-comprehensive-documentation-of-the-turung-and-singpho-languages-of-assam}
\item Stephen Morey's drafts: \url{https://latrobe.academia.edu/StephenMorey/Drafts}
\end{itemize}
\section{Religious and Community Resources}
\begin{itemize}
\item Asia Harvest people-group profile: \url{https://www.people-groups.asiaharvest.org/India/Turung.pdf}
\item Migration history of the Turungs: \url{https://nezine.com/info/SGdDdXZPaGRITTRxazd5RlFlVEFVUT09/migration-history-of-the-turungs,-the-last-tai-group-to-enter-assam.html}
\item Rural India article on cultivating a mother tongue in Assam (the search summary mentions a Turung community documentarian): \url{https://ruralindiaonline.org/article/cultivating-a-mother-tongue-in-assam}
\end{itemize}
\section{Digital and Audio Resources}
\begin{itemize}
\item Endangered Languages Archive collection, Singpho Language of North East India (including Turung): \url{https://eldp.access.preservica.com/dk0006/}
\item Endangered Languages Project entry on an AASU project: \url{https://www.endangeredlanguages.com/node/37647}
\end{itemize}
\section{Curriculum Development Context}
Turung is under pressure from Assamese and English. The ELAR collection and the documentation project are the most substantial sources, and no pedagogical materials were found.
\end{document}